\usepackage[utf8]{inputenc}
\usepackage[russian]{babel}
\usepackage{multirow}
\usepackage[table]{xcolor}
\usepackage{hhline}
\usepackage{array}
\usepackage{tikz}

\usepackage{fontspec}

% Основной шрифт — Times New Roman
\setmainfont{Times New Roman}


\pagestyle{empty}

% Назначение цветов
\definecolor{light-blue}{RGB}{193, 235, 251}
\definecolor{light-yellow}{RGB}{253, 253, 195}
\definecolor{light-red}{RGB}{247, 180, 185}

\definecolor{dark-red}{RGB}{232, 106, 114}
\definecolor{dark-blue}{RGB}{109, 186, 215}
\definecolor{dark-yellow}{RGB}{251, 233, 95}


\newcommand{\placeInCell}[8]{%
    % #1: ширина ячейки, #2: высота ячейки, #3: название элемента слева сверху, #4: расшифровка слева по центру
    % #5: по-русски название слева снизу, #6: порядковый номер справа сверху
    % #7: цвет
    % #8: атомная масса справа посередине
    \begin{tikzpicture}
        \useasboundingbox (0,0) rectangle (#1, #2);
        
        \node[anchor=north west, inner sep=0cm] at (-0.15cm,#2-0.1cm) {\textbf{#3}};
        \node[anchor=north west, inner sep=0cm] at (-0.15cm, #2/2-0.14cm) {{\scriptsize #4}};
        \node[anchor=south west, inner sep=0cm] at (-0.15cm, -0.18cm) {\scriptsize \textbf{#5}};
        \node[anchor=center, inner sep=0.1cm, fill=#7] at (#1-0.35cm, #2-0.18cm) {\scriptsize #6};
        \node[anchor=center, inner sep=0cm] at (#1-0.37cm, #2-0.5cm) {\scriptsize \textbf{#8}};
    \end{tikzpicture}%
}

\newcommand{\placeInCellTwo}[8]{%
    % #1: ширина ячейки, #2: высота ячейки, #3: название элемента справа сверху, #4: расшифровка справа по центру
    % #5: по-русски название справа снизу, #6: порядковый номер слева сверху
    % #7: цвет
    % #8: атомная масса слева посередине
    \begin{tikzpicture}
        \useasboundingbox (0,0) rectangle (#1, #2);
        
        \node[anchor=north east, inner sep=0cm] at (#1+0.15cm,#2-0.1cm) {\textbf{#3}};
        \node[anchor=north east, inner sep=0cm] at (#1+0.15cm, #2/2-0.14cm) {{\scriptsize #4}};
        \node[anchor=south east, inner sep=0cm] at (#1+0.15cm, -0.18cm) {\scriptsize \textbf{#5}};
        \node[anchor=center, inner sep=0.1cm, fill=#7] at (0.35cm, #2-0.18cm) {\scriptsize #6};
        \node[anchor=center, inner sep=0cm] at (0.37cm, #2-0.5cm) {\scriptsize \textbf{#8}};
    \end{tikzpicture}%
}
